\chapter{Introduction}
\label{cha:introduction}

Beer is one of the oldest foods of mankind. Since the Bavarian purity law of 1516 it consists of
only four ingredients: water, malt, hops and yeast \citep{hopfner2019reinheit}. Despite this
simplicity, beers differ enormously in taste and smell -- and a large part of this variety is not
created in the brewhouse but in the fermentation tank \citep{malz2021ester}.

\section{Motivation}
Craft breweries increasingly experiment with fermentation parameters. At the same time there is a
lack of \textbf{systematic data} on how strongly individual parameters actually change the aroma.

\section{Goals and Research Questions}
\begin{enumerate}
  \item How does the concentration of aroma-active esters change with fermentation temperature?
  \item From which temperature do tasters perceive the difference?
  \item Can the aroma profile be predicted with a simple model?
\end{enumerate}

\section{Methodology and Structure}
Chapter~\ref{cha:background} explains the biochemistry of fermentation, Chapter~\ref{cha:design}
the experimental design, Chapter~\ref{cha:evaluation} presents measurements and tasting results.
